\documentclass[10pt]{article}
\usepackage[margin=0.55in]{geometry}
\usepackage[hidelinks]{hyperref}
\usepackage{enumitem}
\usepackage{microtype}
\setlength{\parindent}{0pt}
\setlength{\parskip}{0.35em}
\setlist[itemize]{leftmargin=1.2em, topsep=0.15em, itemsep=0.15em}

\begin{document}

\begin{center}
{\large\textbf{Ethics and Fairness Report}}\\
\textbf{DSCI 601: Applied Data Science I}
\end{center}

My project studies \textbf{fairness-aware sequential decision-making} in two settings:
clinical diagnostic workflows and quantum-network routing. The core question is how
different levels of decision context---from non-contextual to contextual to informative
bandit methods---change fairness, equity, and performance under uncertainty, missing
information, and distribution shift. Although this project is technical, \textbf{fairness} is not
an add-on. It is one of the reasons the project exists at all.

My earlier AI ethics paper for IDAI700 argued that AI systems affecting marginalized
communities should not be evaluated only by average performance or institutional claims of
trustworthiness. They must also be evaluated by who is made visible, whose errors are
ignored, and whether accountability exists when harm is patterned rather than random. That
foundation carries directly into this project. In a clinical workflow, fairness matters because
optimizing overall utility can hide subgroup failures such as false-negative gaps, delayed
follow-up, or weaker decision support for patients whose context is incomplete, delayed, or
noisier. In a quantum-routing workflow, fairness matters because routing and resource
allocation can create service-equity gaps across user or flow groups even when total network
performance looks strong.

Addressing fairness changes the project design and implementation in concrete ways. First,
the system cannot be evaluated on utility alone; it must track disparity over time, including
group-wise outcome gaps rather than only end-of-run averages. Second, the project must
model \textit{context quality} explicitly, because missing or degraded context can itself be a
driver of unfair outcomes. Third, fairness changes the algorithmic comparison: the project
must compare baselines, contextual methods, and informative-context methods not only on
regret or reward, but on whether they reduce or amplify disparities. Finally, fairness makes
mitigation part of implementation rather than an afterthought. At least one fairness-aware
mechanism must be built into the evaluation stack so the work tests whether disparities can be
reduced in practice rather than merely documented.

Fairness also changes what reproducibility and data sharing mean for this project. A
reproducible system should not require unrestricted sharing of sensitive patient data, so the
clinical side is designed simulation-first with controllable missingness, noise, and shift.
That improves reproducibility because others can rerun the same experiments with fixed seeds,
shared configurations, and documented metrics. At the same time, fairness raises the bar for
what must be shared: not only code and aggregate scores, but also metric definitions,
subgroup evaluation procedures, assumptions about groups, and the limitations of the data or
simulation used. In this sense, fairness strengthens reproducibility by forcing the project to
make value-laden design decisions visible instead of hiding them behind average performance.

\vspace{0.2em}
\footnotesize
\textbf{Selected supporting references:} Friedler et al., ``The Inherent Limitations of AI
Fairness''; Givens and Morris, ``Fairness \& Abstraction in Sociotechnical Systems''; Alan
Turing Institute, ``Fairness in AI.'' \\
\textbf{Conceptual foundation:} \textit{IDAI700\_Research\_Paper.pdf} (AI ethics paper, used as
motivation rather than reused text).

\end{document}
